\PassOptionsToPackage{table,xcdraw}{xcolor}
\documentclass[11pt, a4paper]{article}

\usepackage[T1]{fontenc}
\usepackage[utf8]{inputenc}
\usepackage{tgpagella}
\usepackage{tgcursor}
\usepackage[spanish, es-noshorthands]{babel}
\usepackage{amsmath, amssymb}
\usepackage{booktabs}
\usepackage{tabularx}
\usepackage[table]{xcolor}
\usepackage{tikz}
\usetikzlibrary{arrows.meta, positioning, calc, patterns}
\usepackage[margin=1.6cm]{geometry}
\usepackage{enumitem}
\usepackage{microtype}

\emergencystretch=2em
\tolerance=1000
\hbadness=10000

\definecolor{ucbblue}{RGB}{0, 51, 102}
\definecolor{ucbgray}{RGB}{240, 244, 248}
\definecolor{warnred}{RGB}{180, 20, 20}

\newcommand{\workbox}[1]{%
    \par\vspace{0.18cm}\noindent
    \begin{tikzpicture}
        \draw[draw=ucbblue!45, thick, rounded corners=3pt] (0,0) rectangle (\linewidth, #1);
        \foreach \y in {0.7, 1.4, ..., #1} {
            \draw[ucbblue!18, dashed] (0.25, \y) -- (\linewidth-0.25, \y);
        }
        \node[anchor=north west, font=\scriptsize\itshape, text=ucbblue!70] at (0.2, #1-0.08) {Espacio de desarrollo, ecuaciones, unidades e interpretación física:};
    \end{tikzpicture}\par\vspace{0.2cm}%
}

\begin{document}

\noindent
{\LARGE \textbf{\textcolor{ucbblue}{IMT-221 — Circuitos Electrónicos III}}}\\[0.12cm]
{\Large \textbf{Evaluación Escrita Individual EC3 — Cuadernillo para Estudiante (Variante A)}}\\[0.08cm]
{\small \textbf{Sesión confirmada por el docente:} W15S03 — viernes 13 de noviembre de 2026}
\par\vspace{0.15cm}\noindent\rule{\linewidth}{0.6pt}\vspace{0.15cm}

\noindent
\textbf{Nombre:} \rule{6.2cm}{0.4pt} \textbf{Fecha:} 13/11/2026\\[0.2cm]
 \quad \textbf{Turno / Aula:} \rule{4.2cm}{0.4pt} \quad \textbf{Calificación (sobre 40 pt):} \rule{2.4cm}{0.4pt}

\vspace{0.15cm}
\noindent
\fcolorbox{warnred}{ucbgray}{%
\parbox{\dimexpr\linewidth-2\fboxsep-2\fboxrule}{%
\textbf{Duración de resolución:} 75 minutos (dentro de una sesión de 90 min). \quad \textbf{Puntaje interno:} 40 puntos.\\
\textbf{\textcolor{warnred}{Modalidad aprobada: SIN CALCULADORA.}} \quad \textbf{Formulario:} incluido al final de este mismo cuadernillo.\\
No se permite consultar fichas externas ni dispositivos electrónicos; cualquier ajuste institucional debe comunicarse antes del examen.
}}

\vspace{0.15cm}
\noindent
\textbf{Instrucciones:} Puede consultar el formulario incluido al final de este cuadernillo; dibuje e identifique los nodos cuando corresponda; declare modelo e hipótesis; presente ecuación, sustitución, unidades y explicación física. El crédito parcial considera el razonamiento. \textbf{Todos los parámetros ficticios son datos hipotéticos del ejercicio, NO especificaciones comerciales.} No se requiere construir ni usar LTspice.

\par\vspace{0.15cm}\noindent\rule{\linewidth}{0.4pt}

\section*{Problema 1 — Diagnóstico de un interruptor low-side (12 puntos; 20 min)}

Un NMOS ficticio conmuta una carga resistiva equivalente de $\mathbf{9\,\Omega}$ entre el positivo de un bus $\mathbf{10\,\text{V}}$ \textbf{DC} y su drenador D; la fuente S está a GND del bus. Durante ON se ha \textbf{declarado válido}, únicamente para el ejercicio, $R_{DS(\text{on})} = 1\,\Omega$; durante OFF es circuito abierto. Se aplica PWM con $D = 1/2$. Se ignoran transiciones y corrientes de fuga.

\begin{enumerate}[label=\arabic*., noitemsep, topsep=3pt]
    \item (2 pt) Dibuje el lazo $+V_{\text{BUS}} \rightarrow R_L \rightarrow D \rightarrow S \rightarrow \text{GND}$ y localice G, D y S. Indique sentido convencional de corriente.
    \item (4 pt) En ON, escriba KVL y calcule $I_D$, $V_{RL}$ y $V_{DS}$. Verifique que ambas caídas suman el bus.
    \item (3 pt) Calcule las potencias \textbf{en ON} de la carga y del MOSFET. Indique dónde se disipa cada una.
    \item (2 pt) Calcule la potencia \textbf{promedio ideal} de la carga con el PWM indicado y explique la hipótesis temporal necesaria.
    \item (1 pt) En OFF ideal, indique $I_D$, $V_{DS}$ y disipación del MOSFET. ¿La presencia de tensión implica potencia alta?
\end{enumerate}

\section*{Problema 2 — Referencia de gate y validez de una especificación (7 puntos; 12 min)}

Un circuito de prueba tiene $V_G = 10\,\text{V}$ y $V_S = 4\,\text{V}$, ambos referidos a GND. Se comparan dos \textbf{MOSFET ficticios}, cuyos demás requisitos de ensayo se declaran satisfechos:

\begin{center}
\renewcommand{\arraystretch}{1.25}
\begin{tabularx}{0.82\linewidth}{|c|>{\raggedright\arraybackslash}X|}
\hline
\rowcolor{ucbgray}
\textbf{Modelo ficticio} & \textbf{Única garantía proporcionada} \\ \hline
\textbf{X} & $R_{DS(\text{on})} \le 0.20\,\Omega$, especificada con $V_{GS} = 6\,\text{V}$ \\ \hline
\textbf{Y} & $R_{DS(\text{on})} \le 0.10\,\Omega$, especificada con $V_{GS} = 10\,\text{V}$ \\ \hline
\end{tabularx}
\end{center}

\begin{enumerate}[label=\arabic*., noitemsep, topsep=3pt]
    \item (2 pt) Calcule el $V_{GS}$ real. Explique por qué $V_G$ no basta para decidir la excitación.
    \item (3 pt) Determine cuál de las \textbf{dos garantías listadas} puede invocarse bajo el punto de operación y cuál \textbf{no}. ¿Puede afirmarse que Y es necesariamente el de menor pérdida?
    \item (2 pt) El dispositivo Y tiene $V_{GS(\text{th})} = 2\,\text{V}$ \textbf{en una prueba de corriente pequeña}, según la hipótesis. ¿Permite esto afirmar baja resistencia de conducción a $V_{GS} = 6\,\text{V}$? Razone.
\end{enumerate}


\section*{Problema 3 — Carga de gate y energía de un perfil de conmutación (9 puntos; 16 min)}

Se estudia un MOSFET ficticio. La carga de gate transferida en un evento vale $Q_g = 48\,\text{nC}$, y la corriente media de carga durante ese evento se aproxima por $I_G = 12\,\text{mA}$ constante. Para cada encendido, \textbf{la potencia instantánea disipada en el MOSFET} se aproxima por un triángulo de base temporal $4\,\mu\text{s}$ y altura $2\,\text{W}$; para cada apagado, otro triángulo de base $6\,\mu\text{s}$ y altura $1\,\text{W}$. La frecuencia es $f = 1000\,\text{Hz}$. \textit{Los triángulos son una hipótesis explícita, no mediciones.}

\begin{center}
\begin{tikzpicture}[scale=0.88, transform shape]
    % Perfil de encendido (ON)
    \draw[->, thick] (0,0) -- (5.2,0) node[right]{\small $t\,(\mu\text{s})$};
    \draw[->, thick] (0,0) -- (0,2.6) node[above]{\small $p_D(t)\,(\text{W}) \text{ [Hipótesis ON]}$};
    \fill[ucbblue!15] (0.5,0) -- (2.5,2.0) -- (4.5,0) -- cycle;
    \draw[very thick, ucbblue] (0,0) -- (0.5,0) -- (2.5,2.0) -- (4.5,0) -- (5.0,0);
    \draw[dashed] (0,2.0) node[left]{\small $2\,\text{W}$} -- (2.5,2.0);
    \draw[<->, thick] (0.5,-0.35) -- (4.5,-0.35) node[midway, below]{\small $\Delta t_{\text{on}} = 4\,\mu\text{s}$};
    \node[font=\footnotesize\bfseries, text=ucbblue] at (2.5, 0.75) {$E_{\text{on}}$};

    % Perfil de apagado (OFF)
    \begin{scope}[xshift=8.2cm]
        \draw[->, thick] (0,0) -- (6.2,0) node[right]{\small $t\,(\mu\text{s})$};
        \draw[->, thick] (0,0) -- (0,2.6) node[above]{\small $p_D(t)\,(\text{W}) \text{ [Hipótesis OFF]}$};
        \fill[warnred!15] (0.5,0) -- (3.5,1.0) -- (5.5,0) -- cycle;
        \draw[very thick, warnred] (0,0) -- (0.5,0) -- (3.5,1.0) -- (5.5,0) -- (6.0,0);
        \draw[dashed] (0,1.0) node[left]{\small $1\,\text{W}$} -- (3.5,1.0);
        \draw[<->, thick] (0.5,-0.35) -- (5.5,-0.35) node[midway, below]{\small $\Delta t_{\text{off}} = 6\,\mu\text{s}$};
        \node[font=\footnotesize\bfseries, text=warnred] at (3.2, 0.42) {$E_{\text{off}}$};
    \end{scope}
\end{tikzpicture}
\end{center}

\begin{enumerate}[label=\arabic*., noitemsep, topsep=2pt]
    \item (3 pt) A partir de $i_G = dq_G/dt$, estime el tiempo de transferencia de $Q_g$ y conviértalo a $\mu\text{s}$.
    \item (3 pt) Halle $E_{\text{on}}$ y $E_{\text{off}}$ usando el área de los dos triángulos, con unidades de $\mu\text{J}$.
    \item (2 pt) Calcule $P_{\text{sw}} = f(E_{\text{on}} + E_{\text{off}})$ en $\text{mW}$.
    \item (1 pt) Indique \textbf{una razón física} por la que $Q_g/I_G$ no es necesariamente el tiempo exacto observado de caída de $V_{DS}$.
\end{enumerate}

\section*{Problema 4 — PWM, instrumento y referencia de tierra (6 puntos; 12 min)}

Un generador entrega una onda \textbf{cuadrada simétrica} con $V_{\text{pp}} = 6\,\text{V}$ y $\text{offset} = +1\,\text{V}$, ambos valores respecto del borne de retorno del generador. Un estudiante llama a esto «$0\text{ a }6\,\text{V}$». Para observar el drenador de un MOSFET, propone sujetar la \textbf{pinza de masa de un osciloscopio de banco, conectada a tierra de protección}, directamente a D; el GND del bus se encuentra ya unido a tierra de protección por otra conexión.

\begin{enumerate}[label=\arabic*., noitemsep, topsep=3pt]
    \item (2 pt) Calcule $V_{\min}$ y $V_{\max}$ de la señal. ¿La descripción «$0\text{ a }6\,\text{V}$» es válida?
    \item (2 pt) Explique qué conexión peligrosa produciría la pinza de masa del osciloscopio en D y qué puede ocurrir.
    \item (2 pt) Indique \textbf{dos acciones verificables previas al ensayo} para obtener una medición segura y confirmar la excitación real de gate (no basta con escribir «tener cuidado»).
\end{enumerate}

\section*{Problema 5A — Función del driver BJT (6 puntos; 15 min)}

\vspace{0.1cm}
Un driver conceptual presenta una etapa NPN en seguidor de emisor: colector conectado a $+12\,\text{V}$, base comandada por una salida lógica de $\mathbf{3.3\,\text{V}}$ y salida en el emisor. Para este modelo se supone $V_{BE} = 0.7\,\text{V}$ mientras conduce. El NMOS de potencia exige para la condición declarada un nivel de gate de $\mathbf{10\,\text{V}}$ \textbf{respecto a su source} (S a $0\,\text{V}$).

\begin{enumerate}[label=\arabic*., noitemsep, topsep=3pt]
    \item (2 pt) Estime la tensión máxima aproximada de salida del seguidor. Muestre la relación utilizada.
    \item (2 pt) ¿Por qué el hecho de alimentar el colector a $12\,\text{V}$ \textbf{no} garantiza gate a $12\,\text{V}$? ¿Se cumple el objetivo declarado de $10\,\text{V}$?
    \item (2 pt) Distinga \texttt{sourcing} y \texttt{sinking} de carga de gate. Proponga el tipo de función adicional que necesitaría el circuito, sin diseñar ni conectar un driver físico.
\end{enumerate}


\newpage
\section*{FORMULARIO DE ECUACIONES — EC3}

\noindent
\textbf{Integrado en el examen.} Las relaciones siguientes no reemplazan la identificación de las hipótesis, las condiciones de validez ni la interpretación física.

\subsection*{Tensiones, corriente y potencia}

\[V_{GS}=V_G-V_S,\qquad V_{DS}=V_D-V_S.\]

\[V_{\text{BUS}}=I_D R_L+V_{DS},\qquad I_D\approx\frac{V_{\text{BUS}}}{R_L+R_{DS(\text{on})}}\]

\noindent\textit{La aproximación resistiva del MOSFET solo corresponde cuando el problema declara válido el valor de} $R_{DS(\text{on})}$.

\[V=IR,\qquad P=VI=I^2R=\frac{V^2}{R}.\]

\[P_{L}=I_D^2R_L,\qquad P_{\text{MOS,cond}}\approx I_D^2R_{DS(\text{on})}.\]

\subsection*{PWM y transitorios}

\[D=\frac{t_{\text{ON}}}{T},\qquad f=\frac{1}{T},\qquad P_{\text{prom}}\approx D\,P_{\text{ON}}\]

\noindent\textit{La última relación supone potencia despreciable en OFF y transitorios despreciables para el elemento evaluado.}

\[i_G(t)=\frac{dq_G}{dt},\qquad \Delta Q_G=\int i_G(t)\,dt,\qquad t_{\text{transfer}}\approx\frac{Q_G}{I_{G,\text{prom}}}.\]

\[p_D(t)=v_{DS}(t)i_D(t),\qquad E_{\text{evento}}=\int p_D(t)\,dt.\]

\[E_{\text{triangulo}}=\frac{1}{2}\,(\text{base temporal})(\text{altura de potencia}), \qquad P_{\text{sw}}\approx f_{\text{sw}}(E_{\text{on}}+E_{\text{off}}).\]

\subsection*{Señal simétrica y seguidor BJT}

\[V_{\max}=V_{\text{offset}}+\frac{V_{\text{pp}}}{2},\qquad V_{\min}=V_{\text{offset}}-\frac{V_{\text{pp}}}{2}.\]

\[V_{E}\approx V_B-V_{BE}\quad\text{(seguidor NPN, en conducción y bajo las hipótesis del problema).}\]

\subsection*{Unidades útiles}

\[1\,\mathrm{mA}=10^{-3}\,\mathrm{A},\quad 1\,\mathrm{nC}=10^{-9}\,\mathrm{C},\quad 1\,\mu\mathrm{s}=10^{-6}\,\mathrm{s},\quad 1\,\mu\mathrm{J}=10^{-6}\,\mathrm{J}.\]

\vfill
\noindent\rule{\linewidth}{0.4pt}
\begin{center}
\textbf{FIN — Revise unidades, hipótesis y que todas las respuestas sean legibles.}
\end{center}

\end{document}
